\documentclass[11pt, letterpaper]{article}

% --- Idioma y codificación ---
\usepackage[spanish]{babel}
\usepackage[utf8]{inputenc}
\usepackage[T1]{fontenc}
\usepackage{lmodern}

% --- Matemáticas ---
\usepackage{amsmath, amssymb, amsthm}

% --- Formato ---
\usepackage[margin=2.5cm]{geometry}
\usepackage{parskip}
\usepackage{booktabs}
\usepackage{hyperref}
\hypersetup{colorlinks=true, linkcolor=blue!70!black, urlcolor=blue!70!black}

% --- Entornos de teoremas ---
\theoremstyle{definition}
\newtheorem{definition}{Definición}[section]
\theoremstyle{plain}
\newtheorem{theorem}{Teorema}[section]
\newtheorem{lemma}[theorem]{Lema}
\newtheorem{corollary}[theorem]{Corolario}
\theoremstyle{remark}
\newtheorem{remark}{Observación}[section]

\title{%
  T01 --- Demostración formal:\\[4pt]
  Correctitud y complejidad de Dijkstra
}
\author{Curso Linux--C--Rust --- B15 Estructuras de datos en C}
\date{}

\begin{document}
\maketitle
\tableofcontents

% ===================================================================
\section{Modelo}
% ===================================================================

\begin{definition}
Grafo $G = (V, E)$ dirigido con función de peso
$w: E \to \mathbb{R}_{\geq 0}$ (pesos no negativos). Fuente
$s \in V$. Distancia:
$\delta(s, v) = \min\{w(p) : p \text{ es camino de } s \text{ a } v\}$,
donde $w(p) = \sum_{(u,v) \in p} w(u,v)$. Si no existe camino,
$\delta(s,v) = \infty$.
\end{definition}

Dijkstra mantiene: $S$ (vértices finalizados),
$\mathrm{dist}[v]$ (estimación, inicialmente $\infty$;
$\mathrm{dist}[s] = 0$), y cola de prioridad $Q = V \setminus S$.
En cada paso: extraer $u = \arg\min_{v \in Q} \mathrm{dist}[v]$,
agregar $u$ a $S$, relajar aristas salientes de $u$.

% ===================================================================
\section{Lemas previos}
% ===================================================================

\begin{lemma}[Subestructura óptima]\label{lem:substructure}
Si $p = s \leadsto u \to v$ es un camino más corto de $s$ a $v$,
entonces el subcamino $s \leadsto u$ es un camino más corto de $s$
a $u$.
\end{lemma}

\begin{proof}
Si existiera $p'$ de $s$ a $u$ con $w(p') < w(s \leadsto u)$,
entonces $p' \to v$ tendría peso
$w(p') + w(u,v) < w(s \leadsto u) + w(u,v) = w(p)$, contradiciendo
que $p$ es más corto.
\end{proof}

\begin{lemma}[Convergencia de la relajación]\label{lem:relax}
Si $(u,v) \in E$ y $\mathrm{dist}[u] = \delta(s,u)$ al relajar
$(u,v)$, entonces tras la relajación:
$\mathrm{dist}[v] \leq \delta(s,v)$.
\end{lemma}

\begin{proof}
Tras relajar:
$\mathrm{dist}[v] \leq \mathrm{dist}[u] + w(u,v)
= \delta(s,u) + w(u,v)$.
Si el camino más corto a $v$ pasa por $u$:
$\delta(s,v) = \delta(s,u) + w(u,v)$, y la desigualdad vale.
En cualquier caso, $\mathrm{dist}[v]$ nunca baja de $\delta(s,v)$
por el invariante de cota superior (Lema~\ref{lem:upper}).
\end{proof}

\begin{lemma}[Invariante de cota superior]\label{lem:upper}
En todo momento, $\mathrm{dist}[v] \geq \delta(s,v)$ para todo
$v \in V$.
\end{lemma}

\begin{proof}
Por inducción sobre el número de relajaciones.

\emph{Base.} $\mathrm{dist}[s] = 0 = \delta(s,s)$ y
$\mathrm{dist}[v] = \infty \geq \delta(s,v)$ para $v \neq s$.
\checkmark

\emph{Paso.} La relajación de $(u,v)$ pone
$\mathrm{dist}[v] \leftarrow \mathrm{dist}[u] + w(u,v)$.
Por hipótesis inductiva, $\mathrm{dist}[u] \geq \delta(s,u)$:
\[
  \delta(s,v) \leq \delta(s,u) + w(u,v)
  \leq \mathrm{dist}[u] + w(u,v)
  = \mathrm{dist}[v]_{\text{nuevo}} \qedhere
\]
\end{proof}

% ===================================================================
\section{Demostración 1: correctitud}
% ===================================================================

\begin{theorem}\label{thm:correctness}
Al extraer el vértice $u$ de la cola de prioridad,
$\mathrm{dist}[u] = \delta(s,u)$.
\end{theorem}

\begin{proof}
Por contradicción. Sea $u$ el \textbf{primer} vértice extraído con
$\mathrm{dist}[u] \neq \delta(s,u)$. Por el
Lema~\ref{lem:upper}, $\mathrm{dist}[u] > \delta(s,u)$. En
particular, $\delta(s,u) < \infty$ y existe un camino de $s$ a $u$.

\textbf{Paso 1.} Sea $p = s \leadsto u$ un camino más corto. Como
$s \in S$ y $u \notin S$, el camino cruza la frontera de $S$. Sea
$(x,y)$ la primera arista con $x \in S$, $y \notin S$:
\[
  p: \underbrace{s \leadsto x}_{\in S}
  \to \underbrace{y \leadsto u}_{\notin S}
\]

\textbf{Paso 2.} Como $x$ fue extraído antes de $u$ y $u$ es el
primer error, $\mathrm{dist}[x] = \delta(s,x)$. Al extraer $x$, se
relajó $(x,y)$. Por el Lema~\ref{lem:relax}:
\[
  \mathrm{dist}[y] \leq \delta(s,x) + w(x,y) = \delta(s,y)
\]
(igualdad por subestructura óptima, Lema~\ref{lem:substructure}).
Por el Lema~\ref{lem:upper}:
$\mathrm{dist}[y] \geq \delta(s,y)$. Combinando:
$\mathrm{dist}[y] = \delta(s,y)$.

\textbf{Paso 3.} Como $y$ está en el camino óptimo y
$w \geq 0$:
\[
  \mathrm{dist}[y] = \delta(s,y) \leq \delta(s,u)
  < \mathrm{dist}[u]
\]

Pero $u$ fue extraído antes que $y$ (extract-min), así que
$\mathrm{dist}[u] \leq \mathrm{dist}[y]$.
Contradicción.
\end{proof}

% ===================================================================
\section{Demostración 2: necesidad de $w \geq 0$}
% ===================================================================

\begin{theorem}\label{thm:nonneg}
Si existe una arista con peso negativo, Dijkstra puede producir
distancias incorrectas.
\end{theorem}

\begin{proof}
Contraejemplo. Grafo:
$s \xrightarrow{2} a$, $s \xrightarrow{3} b \xrightarrow{-2} a$.
Distancia real: $\delta(s,a) = \min(2, 3{+}({-}2)) = 1$.

Ejecución:
\begin{enumerate}
  \item Extraer $s$: $\mathrm{dist}[a] = 2$,
    $\mathrm{dist}[b] = 3$.
  \item Extraer $a$ ($\mathrm{dist} = 2$). Se finaliza con
    $\mathrm{dist}[a] = 2$.
  \item Extraer $b$ ($\mathrm{dist} = 3$). Relajar $(b,a)$:
    $3 + ({-}2) = 1 < 2$, pero $a$ ya fue finalizado.
\end{enumerate}

Dijkstra retorna $\mathrm{dist}[a] = 2 \neq 1 = \delta(s,a)$.

\textbf{Dónde falla la prueba.} En el Paso~3 del
Teorema~\ref{thm:correctness}, usamos
$\delta(s,y) \leq \delta(s,u)$ porque el subcamino $y \leadsto u$
tiene peso $\geq 0$. Con pesos negativos, podría ser que
$\delta(s,y) > \delta(s,u)$, y la contradicción no se produce.
\end{proof}

% ===================================================================
\section{Demostración 3: complejidad
  \texorpdfstring{$O((n+m)\log n)$}{O((n+m) log n)}}
% ===================================================================

\begin{theorem}\label{thm:complexity}
Dijkstra con min-heap binario ejecuta en $O((n+m)\log n)$.
\end{theorem}

\begin{proof}
Contamos las operaciones sobre el heap:

\emph{Inserciones.} Cada relajación exitosa produce una inserción.
Una arista $(u,v)$ se relaja a lo sumo una vez (al extraer $u$). Hay
a lo sumo $m$ inserciones, cada una $O(\log n)$:
\[
  C_{\text{insert}} = O(m \log n)
\]

\emph{Extracciones.} Se extraen a lo sumo $m$ entradas del heap
(incluyendo obsoletas con lazy deletion), cada una $O(\log n)$:
\[
  C_{\text{extract}} = O(m \log n)
\]

Total, incluyendo inicialización $O(n)$:
\[
  T = O(n + m\log n) = O((n+m)\log n) \qedhere
\]
\end{proof}

% ===================================================================
\section{Demostración 4: optimalidad del array para grafos densos}
% ===================================================================

\begin{theorem}\label{thm:array}
Dijkstra con array lineal ejecuta en $O(n^2)$, que es óptimo para
grafos con $m = \Omega(n^2/\log n)$.
\end{theorem}

\begin{proof}
Con array lineal: \texttt{extract\_min} escanea $n$ elementos en
$O(n)$, ejecutado $n$ veces: $O(n^2)$. Relajación: $O(1)$ por
arista, total $O(m)$.
\[
  T_{\text{array}} = O(n^2 + m) = O(n^2)
  \quad \text{(si } m = O(n^2)\text{)}
\]

Con heap: $T_{\text{heap}} = O((n+m)\log n)$. Para
$m = \Theta(n^2)$:
\[
  T_{\text{heap}} = O(n^2 \log n) > O(n^2) = T_{\text{array}}
\]

Punto de cruce: $m = n^2/\log n$. Para grafos más densos, el
array gana.
\end{proof}

% ===================================================================
\section{Demostración 5: los arcos parent forman un árbol}
% ===================================================================

\begin{theorem}\label{thm:tree}
Los arcos
$\{(\mathrm{parent}[v], v) : v \neq s,\, \delta(s,v) < \infty\}$
forman un árbol enraizado en~$s$.
\end{theorem}

\begin{proof}
Sea $T_\pi$ el subgrafo inducido por los arcos parent. Sea $V'$ el
conjunto de vértices alcanzables.

\textbf{Conexo desde $s$.} Para todo $v$ con
$\delta(s,v) < \infty$, la cadena
$v, \mathrm{parent}[v], \mathrm{parent}[\mathrm{parent}[v]], \ldots$
termina en $s$ (cada parent tiene distancia estrictamente menor, y
la distancia es finita y no negativa).

\textbf{Acíclico.} Si hubiera un ciclo
$v_1 \to v_2 \to \cdots \to v_k \to v_1$, entonces
$\mathrm{dist}[v_1] > \mathrm{dist}[v_2] > \cdots >
\mathrm{dist}[v_k] > \mathrm{dist}[v_1]$, lo que es imposible.

\textbf{$|V'| - 1$ aristas.} Cada $v \neq s$ tiene exactamente un
parent: $|E(T_\pi)| = |V'| - 1$.

Conexo, acíclico, $|V'| - 1$ aristas: $T_\pi$ es un árbol.
\end{proof}

% ===================================================================
\section{Resumen de resultados}
% ===================================================================

\begin{table}[h]
\centering
\begin{tabular}{llp{5.5cm}}
\toprule
\textbf{Resultado} & \textbf{Técnica} & \textbf{Clave} \\
\midrule
$\mathrm{dist}[u] = \delta(s,u)$ al extraer
  & Contradicción
  & Primer vértice fuera de $S$ en camino óptimo tiene
    $\mathrm{dist}[y] \leq \mathrm{dist}[u]$ \\
Necesidad de $w \geq 0$
  & Contraejemplo
  & Vértice finalizado prematuramente \\
$O((n{+}m)\log n)$ con heap
  & Contar operaciones
  & $\leq m$ inserciones y extracciones, $O(\log n)$ c/u \\
$O(n^2)$ óptimo si $m = \Omega(n^2/\log n)$
  & Comparación directa
  & $n^2 < m\log n$ cuando $m > n^2/\log n$ \\
Arcos parent: árbol
  & Conexo + acíclico + conteo
  & Cadena finita; sin ciclos por dist.\ decrecientes \\
\bottomrule
\end{tabular}
\end{table}

\end{document}
